\documentclass[11pt]{article}

\usepackage[margin=1in]{geometry}
\usepackage{amsmath,amssymb,amsthm}
\usepackage{mathtools}
\usepackage{stmaryrd}
\usepackage{enumitem}
\usepackage{booktabs}
\usepackage{array}
\usepackage{hyperref}
\usepackage{xcolor}

\hypersetup{
  colorlinks=true,
  linkcolor=blue!70!black,
  citecolor=blue!70!black,
  urlcolor=blue!70!black
}

\newtheorem{theorem}{Theorem}[section]
\newtheorem{lemma}[theorem]{Lemma}
\newtheorem{proposition}[theorem]{Proposition}
\newtheorem{corollary}[theorem]{Corollary}
\newtheorem{definition}[theorem]{Definition}
\newtheorem{remark}[theorem]{Remark}

\newcommand{\Code}{\mathsf{Code}}
\newcommand{\CExp}{\mathsf{CExp}}
\newcommand{\CVar}{\mathsf{CVar}}
\newcommand{\Formula}{\mathsf{Formula}}
\newcommand{\Term}{\mathsf{Term}}
\newcommand{\Var}{\mathsf{Var}}
\newcommand{\Proof}{\mathsf{Proof}}
\newcommand{\ProofN}{\mathsf{ProofN}}
\newcommand{\ProofExt}{\mathsf{ProofExt}}
\newcommand{\Schema}{\mathsf{Schema}}
\newcommand{\Nat}{\mathbb{N}}
\newcommand{\enc}[1]{\ulcorner #1 \urcorner}
\newcommand{\evalCExp}{\mathsf{evalCExp}}
\newcommand{\evalCExpN}{\mathsf{evalCExpN}}
\newcommand{\checkProof}{\mathsf{checkProof}}
\newcommand{\checkProofN}{\mathsf{checkProofN}}
\newcommand{\decFormula}{\mathsf{decFormula}}
\newcommand{\encFormula}{\mathsf{encFormula}}
\newcommand{\substF}{\mathsf{substFormulaCode0}}
\newcommand{\csub}{\mathsf{csub}}
\newcommand{\clit}{\mathsf{clit}}
\newcommand{\ccheck}{\mathsf{ccheck}}
\newcommand{\cvar}{\mathsf{cvar}}
\newcommand{\catom}{\mathsf{catom}}
\newcommand{\cnode}{\mathsf{cnode}}
\newcommand{\fbot}{\bot}
\newcommand{\feq}{\mathsf{feq}}
\newcommand{\fimp}{\mathsf{fimp}}
\newcommand{\fall}{\mathsf{fall}}
\newcommand{\fcode}{\mathsf{fcode}}
\newcommand{\fceq}{\mathsf{fceq}}
\newcommand{\fcAll}{\mathsf{fcAll}}
\newcommand{\TrueF}{\mathsf{TrueF}}
\newcommand{\TrueFormula}{\mathsf{TrueFormula}}
\newcommand{\MetaNot}{\mathsf{MetaNot}}
\newcommand{\GS}{\mathcal{G}}
\newcommand{\phiCode}{\mathsf{phiCode}}
\newcommand{\Set}{\mathsf{Set}}

\title{Syntax-Native Incompleteness:\\
  A Chwistek-Inspired Formalization in Agda}
\author{}
\date{}

\begin{document}

\maketitle

\begin{abstract}
We present a fully constructive formalization of G\"odel's first incompleteness theorem
in Agda (3578 lines, 16 files) that uses binary trees as the native code representation,
avoiding the traditional detour through arithmetic and G\"odel numbering entirely.
Self-reference is achieved via a finite quine construction in the style of Chwistek's
metasystem/object-system separation. We develop two architectures---\emph{stratified}
and \emph{fuel-based}---and prove G\"odel~I unconditionally in the stratified setting.
A reflection hierarchy theorem shows that the stratified architecture \emph{necessarily}
separates meta-levels: each proof system can reflect the layer below but is provably
blind to its own reflective reasoning. The fuel-based architecture collapses this
hierarchy, enabling all three Hilbert--Bernays derivability conditions and a
meta-level bounded G\"odel~II. The entire development uses no postulates, no standard
library, no holes, and only ASCII identifiers; the formalization is self-contained and
compiles with Agda's \texttt{--without-K --exact-split} flags.
\end{abstract}

%% ====================================================================
\section{Introduction}
\label{sec:intro}

G\"odel's incompleteness theorems are among the most celebrated results in
mathematical logic, yet their formalizations typically follow a well-worn path:
encode syntax as natural numbers via a G\"odel numbering, develop significant
arithmetic machinery (the $\beta$-function lemma, $\Sigma_1$-completeness),
and then prove that the proof predicate is representable. This path is
mathematically natural but computationally heavy and historically tied to
the specific framework of Peano arithmetic.

Chwistek's approach to metamathematics~\cite{Chwistek1939} offers an
alternative perspective. Rather than embedding syntax into arithmetic,
Chwistek maintained a strict separation between the \emph{object system}
and the \emph{metasystem}, manipulating expressions directly as syntactic
objects. Guard~\cite{Guard1969} later explored a single-system approach with
an internal proof predicate, defining $D_f(x) = Sb(x, n(x))$ to achieve
self-reference by applying a formula to its own number.

In this paper, we pursue a \emph{syntax-native} approach inspired by
Chwistek's philosophy: syntax is represented as binary trees (type $\Code$),
not as natural numbers. Self-reference is achieved through a finite quine
construction ($\csub$) that decodes a formula from its tree code, substitutes,
and re-encodes. The proof system is Hilbert-style with a decidable proof
checker operating directly on tree codes. No arithmetic is needed.

\medskip
\noindent\textbf{Contributions.}
\begin{enumerate}[label=(\arabic*),nosep]
  \item A complete Agda formalization of G\"odel~I with \emph{no assumptions}
    (no consistency hypothesis, no soundness postulate) in the stratified
    architecture (\S\ref{sec:goedel1}).
  \item A \emph{reflection hierarchy theorem} showing that the stratified
    evaluator is provably blind to extended proof codes, making the
    meta-level separation structurally necessary (\S\ref{sec:hierarchy}).
  \item A fuel-based architecture where all three Hilbert--Bernays
    derivability conditions hold and a bounded G\"odel~II is proved at
    the meta-level (\S\ref{sec:fuel}, \S\ref{sec:goedel2}).
  \item A sharp comparison of the two architectures, connecting them to
    Chwistek's stratified metamathematics and Guard's single-system
    approach (\S\ref{sec:comparison}).
\end{enumerate}

%% ====================================================================
\section{The Object Language}
\label{sec:language}

\subsection{Code: Binary Trees as Syntax}

The representation medium is the type of binary trees with natural-number
leaves:
\[
  \Code ::= \catom(n) \mid \cnode(c_1, c_2)
  \qquad (n \in \Nat)
\]
Every syntactic category---variables, terms, formulas, schemas, proof
trees---is encoded into $\Code$ via a tag scheme. For example,
$\catom(3)$ encodes $\fbot$, $\cnode(\catom(5), \cnode(\enc{A}, \enc{B}))$
encodes $A \to B$, and $\cnode(\catom(30), \enc{t})$ encodes the
\texttt{ax-refl} proof of $t = t$. Tags 0--9 are reserved for formulas,
10--13 for schemas, 14--19 for code expressions, and 30--36 for proof
rules.

\subsection{Terms and Formulas}

Terms are built from de~Bruijn--indexed variables, zero, and successor:
\[
  \Term ::= \mathsf{tvar}(v) \mid \mathsf{tzero} \mid \mathsf{tsucc}(t)
\]
Formulas extend a minimal first-order language with code-level connectives:
\begin{align*}
  \Formula ::=\;
  & \fbot
    \mid \feq(s, t)
    \mid \fimp(A, B)
    \mid \fall(A) \\
  &\mid \fcode(c)
    \mid \fceq(e_1, e_2)
    \mid \fcAll(A)
\end{align*}
Here $\fcode(c)$ embeds a code value as a formula, $\fceq(e_1, e_2)$
asserts that two code expressions evaluate to the same code, and
$\fcAll(A)$ universally quantifies over code values (binding a code
variable, distinct from the term variables bound by $\fall$).

Code expressions are the computational fragment:
\[
  \CExp ::= \cvar(v) \mid \clit(c) \mid \ccheck(e) \mid \csub(e_1, e_2)
\]
The expression $\clit(c)$ is a literal code value; $\ccheck(e)$ evaluates
$e$ to a proof code $p$, runs the proof checker, and returns the encoding
of the proved formula; $\csub(e_1, e_2)$ decodes $e_1$ as a formula,
substitutes $\clit(e_2)$ for code variable~0, and re-encodes. This last
operation is the engine of the quine trick.

\subsection{Proof System}

The proof system is Hilbert-style with six constructors:
\begin{align*}
  &\mathsf{ax\text{-}refl}(t) : \Proof(\feq(t, t)) \\
  &\mathsf{ax\text{-}k}(A, B) : \Proof(A \to (B \to A)) \\
  &\mathsf{ax\text{-}s}(A, B, C) : \Proof((A \to B \to C) \to (A \to B) \to A \to C) \\
  &\mathsf{mp}(p_1, p_2) : \Proof(A \to B) \to \Proof(A) \to \Proof(B) \\
  &\mathsf{gen}(p) : \Proof(A) \to \Proof(\fall A) \\
  &\mathsf{cgen}(p) : \Proof(A) \to \Proof(\fcAll A)
\end{align*}
The checker function $\checkProof : \Code \to \mathsf{Maybe}\;\Formula$
dispatches on tag numbers: tag~30 for \texttt{ax-refl}, 31 for
\texttt{ax-k}, 32 for \texttt{ax-s}, 33 for \texttt{mp} (which
recursively checks both subproofs), 34 for \texttt{gen}, and 35 for
\texttt{cgen}. It is structurally recursive on $\Code$ and total.

\subsection{Encoding and Decoding}

The encoding functions $\encFormula : \Formula \to \Code$,
$\mathsf{encTerm} : \Term \to \Code$, etc., are straightforward
recursive mappings. The corresponding decoders $\decFormula$,
$\mathsf{decTerm}$, etc., are partial (returning $\mathsf{Maybe}$).
The key roundtrip lemma is:

\begin{lemma}[Roundtrip]
  \label{lem:roundtrip}
  For all $A : \Formula$,
  $\decFormula(\encFormula(A)) = \mathsf{just}(A)$.
\end{lemma}

Analogous roundtrip lemmas hold for terms, code variables, code
expressions, and schemas.

%% ====================================================================
\section{Self-Reference via the Quine Trick}
\label{sec:quine}

\subsection{Schemas and Diagonalization}

A \emph{schema} is a formula template with a single hole, to be filled
by a code value:
\[
  \Schema ::= \mathsf{shole}
    \mid \mathsf{sconst}(A)
    \mid \mathsf{simp}(S_1, S_2)
    \mid \mathsf{sall}(S)
\]
Instantiation $\mathsf{instantiate}(S, c)$ fills the hole with
$\fcode(c)$, and encoding $\mathsf{encSchema}$ maps schemas to $\Code$.

The \emph{diagonal} of a schema feeds it its own code:
\[
  \mathsf{diag}(S) = \mathsf{instantiate}(S, \mathsf{encSchema}(S))
\]
This is finite, terminating, and definitionally equal to its unfolding.
It is the syntax-native analogue of Guard's $D_f(x) = Sb(x, n(x))$.

\subsection{The G\"odel Sentence}

The G\"odel sentence is constructed in three steps.

\medskip
\noindent\textbf{Step 1: The template $\varphi$.}
Define $\varphi$ as a formula with one free code variable (index~0):
\[
  \varphi = \fcAll\bigl(\fimp\bigl(
    \fceq(\ccheck(\cvar_1),\; \csub(\cvar_0, \cvar_0)),\;
    \fbot\bigr)\bigr)
\]
Here $\cvar_0$ will receive the self-code, and $\cvar_1$ is bound by
$\fcAll$ (it ranges over proof codes). Intuitively, $\varphi$ says:
``for all proof codes $p$, what $p$ proves is not what $\csub$ computes
from my code.''

\medskip
\noindent\textbf{Step 2: The self-code.}
Let $\phiCode = \encFormula(\varphi) : \Code$.

\medskip
\noindent\textbf{Step 3: The G\"odel sentence.}
\[
  \GS = \substF\bigl(\clit(\phiCode),\; \varphi\bigr)
\]
Substituting $\clit(\phiCode)$ for code variable~0 in $\varphi$ yields:
\[
  \GS = \fcAll\bigl(\fimp\bigl(
    \fceq(\ccheck(\cvar_0),\; \csub(\clit(\phiCode), \clit(\phiCode))),\;
    \fbot\bigr)\bigr)
\]

\begin{theorem}[Self-Reference]
  \label{thm:selfref}
  $\evalCExp(\csub(\clit(\phiCode), \clit(\phiCode)))
    = \mathsf{just}(\encFormula(\GS))$.
\end{theorem}

\begin{proof}
  By computation. The $\csub$ evaluator retrieves $\phiCode$, decodes it
  to $\varphi$ (by the roundtrip lemma), substitutes $\clit(\phiCode)$
  for code variable~0, and re-encodes. The result is
  $\encFormula(\substF(\clit(\phiCode), \varphi)) = \encFormula(\GS)$.
\end{proof}

Thus $\GS$ says: ``for all proof codes $p$, $\checkProof(p)$ does not
equal $\encFormula(\GS)$''---i.e., ``I am not provable.'' Crucially,
$\GS$ does \emph{not} contain $\encFormula(\GS)$ as a subexpression
(which would require an infinite term). It contains the finite expression
$\csub(\clit(\phiCode), \clit(\phiCode))$, which \emph{computes}
$\encFormula(\GS)$ at evaluation time.

%% ====================================================================
\section{G\"odel's First Incompleteness Theorem (Stratified)}
\label{sec:goedel1}

\subsection{Soundness}

Truth is defined via an environment-based predicate
$\TrueF(\rho, \sigma, A)$, structurally recursive on $\Formula$:
\begin{itemize}[nosep]
  \item $\TrueF(\rho, \sigma, \fbot) = \varnothing$
  \item $\TrueF(\rho, \sigma, \feq(s,t)) = (\mathsf{eval}_\rho(s) \equiv \mathsf{eval}_\rho(t))$
  \item $\TrueF(\rho, \sigma, A \to B) = \TrueF(\rho,\sigma,A) \to \TrueF(\rho,\sigma,B)$
  \item $\TrueF(\rho, \sigma, \fall A) = \forall t.\; \TrueF(\rho[0 \mapsto t], \sigma, A)$
  \item $\TrueF(\rho, \sigma, \fceq(e_1,e_2)) = \exists c.\; \mathsf{eval}_\sigma(e_1) = \mathsf{just}(c) \,\wedge\, \mathsf{eval}_\sigma(e_2) = \mathsf{just}(c)$
  \item $\TrueF(\rho, \sigma, \fcAll A) = \forall c.\; \TrueF(\rho, \sigma[0 \mapsto c], A)$
\end{itemize}
A formula $A$ is \emph{true}, written $\TrueFormula(A)$, if $\TrueF(\rho, \sigma, A)$
holds for all environments $\rho, \sigma$.

\begin{theorem}[Soundness]
  \label{thm:soundness}
  If $\Proof(A)$, then $\TrueFormula(A)$.
\end{theorem}

\begin{proof}
  By induction on $\Proof(A)$. Each case is direct:
  \texttt{ax-refl} gives $\mathsf{refl}$;
  \texttt{ax-k} gives $\lambda a.\, \lambda \_.\, a$;
  \texttt{ax-s} gives $\lambda f.\, \lambda g.\, \lambda a.\, f\,a\,(g\,a)$;
  \texttt{mp} applies the induction hypotheses;
  \texttt{gen} and \texttt{cgen} abstract over the new bound variable.
\end{proof}

\subsection{Proof Encoding}

Every $\Proof$ tree is encoded as a $\Code$ value by
$\mathsf{encodeProof}$, and the checker accepts these codes:

\begin{lemma}[Encoding Correctness]
  \label{lem:encode}
  For all $\mathit{prf} : \Proof(A)$,
  $\checkProof(\mathsf{encodeProof}(\mathit{prf})) = \mathsf{just}(A)$.
\end{lemma}

\subsection{The Main Theorem}

\begin{theorem}[G\"odel I, No Assumptions]
  \label{thm:goedel1}
  $\MetaNot(\Proof(\GS))$, i.e., there is no proof of the G\"odel sentence.
\end{theorem}

\begin{proof}
  Suppose $\mathit{prf} : \Proof(\GS)$.
  \begin{enumerate}[nosep]
    \item By Theorem~\ref{thm:soundness},
      $\TrueF(\rho_\varnothing, \sigma_\varnothing, \GS)$ holds.
    \item Let $p = \mathsf{encodeProof}(\mathit{prf})$. Since
      $\GS = \fcAll(\ldots)$, instantiate the universal at $c = p$:
      \[
        \TrueF(\rho_\varnothing, \sigma_\varnothing[0 \mapsto p], \mathsf{GoedelBody})
      \]
      where $\mathsf{GoedelBody} = \fimp(\fceq(\ccheck(\cvar_0), \csub(\clit(\phiCode), \clit(\phiCode))), \fbot)$.
      This means: if the two code expressions evaluate to the same value
      under $\sigma_\varnothing[0 \mapsto p]$, then $\varnothing$ (contradiction).
    \item By Lemma~\ref{lem:encode},
      $\checkProof(p) = \mathsf{just}(\GS)$, so
      \[
        \mathsf{eval}_{\sigma[0\mapsto p]}(\ccheck(\cvar_0))
          = \mathsf{just}(\encFormula(\GS)).
      \]
    \item By Theorem~\ref{thm:selfref},
      \[
        \mathsf{eval}_{\sigma[0\mapsto p]}(\csub(\clit(\phiCode), \clit(\phiCode)))
          = \mathsf{just}(\encFormula(\GS)).
      \]
    \item Both sides evaluate to $\encFormula(\GS)$, so the code-equality
      witness exists. Applying step~2 yields $\varnothing$: contradiction.
      \qedhere
  \end{enumerate}
\end{proof}

\begin{remark}
  No consistency hypothesis is assumed. The proof derives a
  contradiction in Agda's metatheory from the bare existence of a
  $\Proof(\GS)$ term. This is possible because soundness is proved
  \emph{within} Agda for the specific (weak) object proof system,
  and the encoding correctness is verified computationally.
\end{remark}

%% ====================================================================
\section{The Reflection Hierarchy}
\label{sec:hierarchy}

\subsection{The Extended Proof System}

To approach G\"odel~II, we extend the proof system with a reflection rule:
\[
  \mathsf{ax\text{-}eval}(e, c, \mathit{eq}) :
    \ProofExt(\fceq(e, \clit(c)))
    \qquad\text{when } \evalCExp(e) = \mathsf{just}(c)
\]
This rule internalizes evaluation results: if a closed code expression
$e$ evaluates to $c$, the extended system can prove the corresponding
$\fceq$ formula. The extended checker $\mathsf{checkProofExt}$ handles
all tags 30--36, where tag~36 is \texttt{ax-eval}.

The first derivability condition holds for base proofs:

\begin{proposition}[D1 for Base Proofs]
  For any $\mathit{prf} : \Proof(A)$,
  \[
    \ProofExt\bigl(\fceq(\ccheck(\clit(\mathsf{encodeProof}(\mathit{prf}))),\;
      \clit(\encFormula(A)))\bigr).
  \]
\end{proposition}

\subsection{The Blindness Lemma}

The stratified architecture creates an asymmetry: the evaluator
$\evalCExp$ calls the \emph{old} checker $\checkProof$ (tags 30--35),
not the extended $\mathsf{checkProofExt}$ (tags 30--36). Consequently:

\begin{lemma}[Blindness]
  \label{lem:blind}
  For any \texttt{ax-eval} proof $\mathit{prf} = \mathsf{ax\text{-}eval}(e, c, \mathit{eq})$,
  \[
    \evalCExp(\ccheck(\clit(\mathsf{encodeProofExt}(\mathit{prf})))) = \mathsf{nothing}.
  \]
\end{lemma}

\begin{proof}
  The encoding $\mathsf{encodeProofExt}(\mathsf{ax\text{-}eval}(e, c, \mathit{eq}))$
  has tag~36. The old $\checkProof$ dispatches through tags 30--35 and
  falls through to $\mathsf{nothing}$. This holds \emph{definitionally}
  (by \texttt{refl} in Agda).
\end{proof}

\begin{theorem}[No Self-Reflection]
  \label{thm:noselfref}
  For any \texttt{ax-eval} proof $\mathit{prf}$,
  there is no code $d$ such that
  $\evalCExp(\ccheck(\clit(\mathsf{encodeProofExt}(\mathit{prf})))) = \mathsf{just}(d)$.
\end{theorem}

\begin{proof}
  By the blindness lemma, the left-hand side is $\mathsf{nothing}$,
  and $\mathsf{nothing} \neq \mathsf{just}(d)$ for any $d$.
\end{proof}

\subsection{Strictness of the Hierarchy}

Combining these results:
\begin{itemize}[nosep]
  \item \textbf{Positive:} For any old $\Proof(A)$, the extended system
    can internally state that the proof code checks to~$A$.
  \item \textbf{Negative:} For any \texttt{ax-eval} proof, the extended
    system \emph{cannot} state the same thing using \texttt{ax-eval},
    because the old evaluator is blind to tag~36.
\end{itemize}

The pattern repeats: a second extension reflecting
$\mathsf{checkProofExt}$ would be blind to its own tag, requiring a
third extension, and so on. Each consistency or completeness result
requires a strictly stronger metasystem---exactly as Chwistek's
framework predicts.

\begin{corollary}
  D3 (the third Hilbert--Bernays derivability condition) is
  \emph{impossible} in the stratified architecture. Consequently,
  G\"odel~II cannot be proved in this setting.
\end{corollary}

%% ====================================================================
\section{The Fuel-Based Architecture}
\label{sec:fuel}

\subsection{Unified Mutual Recursion}

To overcome the D3 barrier, we replace the stratified evaluator/checker
pair with a \emph{fuel-based} mutual recursion:
\begin{align*}
  &\checkProofN : \Nat \to \Code \to \mathsf{Maybe}\;\Formula \\
  &\evalCExpN : \Nat \to \CExp \to \mathsf{Maybe}\;\Code
\end{align*}
Both functions decrease the fuel parameter at each mutual call, ensuring
termination. Crucially, $\evalCExpN$'s $\ccheck$ case calls
$\checkProofN$ (which handles tag~36), not the old $\checkProof$.
Tag~36 in $\checkProofN$ calls $\evalCExpN$ to verify the evaluation
result at runtime, guarded by a Boolean equality check on the code.

\subsection{All Three Derivability Conditions}

\begin{theorem}[D1, Fuel-Based]
  \label{thm:d1fuel}
  Given $\checkProofN(n{+}1, p) = \mathsf{just}(A)$, there exists a
  code $q$ with
  \[
    \checkProofN(n{+}4, q) =
      \mathsf{just}\bigl(\fceq(\ccheck(\clit(p)),\, \clit(\encFormula(A)))\bigr).
  \]
  Each level of representation costs 3 additional fuel.
\end{theorem}

\begin{proof}
  The witness $q = \cnode(\catom(36), \cnode(\mathsf{encCExp}(\ccheck(\clit(p))), \encFormula(A)))$.
  The proof chains through: (1)~decoding the $\CExp$ encoding (by the
  roundtrip lemma), (2)~verifying that $\evalCExpN$ at the appropriate
  fuel yields $\encFormula(A)$ (by the hypothesis), and
  (3)~reflexivity of code equality.
\end{proof}

D2 (closure under modus ponens at the code level) follows by the same
tag-33 construction as in the base system. D3 is simply D1 applied twice:

\begin{theorem}[D3, Fuel-Based]
  \label{thm:d3fuel}
  Given $\checkProofN(n{+}1, p) = \mathsf{just}(A)$, there exists a
  code $q'$ accepted at fuel $n{+}6$ as a proof that the D1
  representation is itself represented.
\end{theorem}

\begin{proof}
  Apply Theorem~\ref{thm:d1fuel} to the D1 witness itself.
\end{proof}

The hierarchy has \emph{collapsed}: unlike the stratified architecture,
where the evaluator was blind to tag~36, the fuel-based $\checkProofN$
handles tag~36 at every fuel level. Self-reflection works because the
checker sees its own reflected proofs.

%% ====================================================================
\section{Bounded G\"odel II}
\label{sec:goedel2}

\subsection{Fuel-Based Soundness}

The fuel-based proof type $\ProofN(n)$ parameterizes proofs by fuel
level. It includes all base constructors plus $\mathsf{axEvalN}(e, c, \mathit{eq})$
when $\evalCExpN(n, e) = \mathsf{just}(c)$.

\begin{theorem}[Soundness for $\ProofN$]
  \label{thm:soundN}
  If $\ProofN(n, A)$, then $\mathsf{TrueFormulaN}(n, A)$, where
  $\mathsf{TrueFormulaN}$ is the truth predicate using $\evalCExpN$ at
  fuel~$n$.
\end{theorem}

\subsection{Bounded G\"odel I}

\begin{theorem}[Bounded G\"odel I]
  \label{thm:goedel1fuel}
  If\/ $\mathit{prf} : \ProofN(n{+}3, \GS)$ and
  $\checkProofN(n{+}2, \mathsf{encodeProofN}(\mathit{prf}))
    = \mathsf{just}(\GS)$,
  then $\varnothing$ (Empty).
\end{theorem}

\begin{proof}
  The same structure as Theorem~\ref{thm:goedel1}: soundness gives
  truth of $\GS$, the encoding gives a proof code, and the self-reference
  property (which holds for fuel $\geq 2$ since $\csub$ of literals
  does not invoke the checker) produces the contradiction.
\end{proof}

The encoding correctness assumption
($\checkProofN(n{+}2, \mathsf{encodeProofN}(\mathit{prf})) = \mathsf{just}(\GS)$)
is the fuel-based analogue of Lemma~\ref{lem:encode}. In the stratified
setting, this was proved by induction on the proof tree. In the
fuel-based setting, the same proof would go through but requires
adapting the entire encoding correctness chain, which is left as the
\texttt{enc-correct} parameter.

\subsection{Meta-Level Bounded G\"odel II}

Bounded consistency is expressed as: no proof code checked at fuel~$n$
proves $\fbot$. The consistency formula $\mathsf{Con}$ is:
\[
  \mathsf{Con} = \fcAll\bigl(\fimp\bigl(
    \fceq(\ccheck(\cvar_0),\; \clit(\encFormula(\fbot))),\;
    \fbot\bigr)\bigr)
\]

\begin{theorem}[Meta-Level Bounded G\"odel II]
  \label{thm:goedel2}
  If\/ $\ProofN(n{+}3, \GS)$ with encoding correctness, then
  $\varnothing$.
  Hence $\ProofN(n{+}3, \mathsf{Con})$ and
  $\ProofN(n{+}3, \GS)$ (with encoding correctness) cannot
  both exist.
\end{theorem}

The full \emph{internal} G\"odel~II---deriving
$\ProofN(m, \fimp(\mathsf{Con}, \GS))$---requires the system to prove
its own soundness internally, which is the classical barrier. The
fuel-based architecture provides all the necessary derivability
conditions (D1, D2, D3), but the final step of internalizing soundness
remains a universal limitation, not an architectural one.

%% ====================================================================
\section{Comparison: Chwistek, Guard, and the Two Architectures}
\label{sec:comparison}

\subsection{Chwistek's Stratified Metamathematics}

Chwistek~\cite{Chwistek1939} insisted on a strict separation between
object system and metasystem, manipulating expressions directly rather
than through arithmetic encoding. Our stratified architecture embodies
this philosophy: $\evalCExp$ calls $\checkProof$ (the object-level
checker), and $\mathsf{checkProofExt}$ lives at a higher meta-level.
The blindness lemma (Lemma~\ref{lem:blind}) \emph{proves} that this
separation is not merely presentational but structurally necessary.

\subsection{Guard's Single-System Approach}

Guard~\cite{Guard1969} pursued a single-system approach with an internal
proof predicate, defining self-reference via $D_f(x) = Sb(x, n(x))$.
Our fuel-based architecture corresponds to this: $\checkProofN$ and
$\evalCExpN$ form a single mutual recursion where the checker can see
its own reflected proofs. The derivability conditions all hold, but at
the cost of a fuel parameter that bounds the depth of self-reflection.

\subsection{Sharp Comparison}

\begin{table}[h]
\centering
\begin{tabular}{@{}lcc@{}}
\toprule
& \textbf{Stratified} & \textbf{Fuel-Based} \\
\midrule
G\"odel I & Proved (no assumptions) & Proved (with \texttt{enc-correct}) \\
D1 & Yes (base proofs only) & Yes (all proofs) \\
D3 & Blocked (blind to tag 36) & Works (+3 fuel per level) \\
Hierarchy & Strict (proved) & Collapses \\
G\"odel II & Impossible & Meta-level proved \\
\bottomrule
\end{tabular}
\caption{Comparison of the two architectures.}
\label{tab:comparison}
\end{table}

The stratified architecture gives a \emph{stronger} G\"odel~I (no
assumptions at all) and a richer structural theorem (the hierarchy), but
cannot reach G\"odel~II. The fuel-based architecture gives weaker
individual theorems (encoding correctness is assumed) but provides the
full derivability infrastructure needed for G\"odel~II.

%% ====================================================================
\section{Related Work}
\label{sec:related}

Formalizations of the incompleteness theorems in proof assistants include
the work of Shankar~\cite{Shankar1994} in Nqthm, O'Connor's
formalization in Coq~\cite{OConnor2005}, and Paulson's formalization
of the incompleteness theorems in Isabelle/HOL~\cite{Paulson2015}. All
of these follow the arithmetic route, encoding syntax into natural
numbers and developing substantial number-theoretic infrastructure.

Our approach is closer in spirit to the ``syntax-based'' formalizations
explored by Rathjen and Leigh~\cite{Leigh2010} and to the category of
``direct'' or ``semantic'' proofs of incompleteness, where the diagonal
lemma is applied at the level of syntax rather than arithmetic. The key
novelty is the use of binary trees as the native representation, which
eliminates the need for G\"odel numbering entirely and makes the
self-reference construction (via $\csub$) both finite and transparent.

The reflection hierarchy theorem connects to the well-known phenomenon
of ordinal analysis and proof-theoretic strength: each consistency
statement requires a strictly stronger system. Our formalization makes
this structural necessity explicit at the level of a concrete proof
checker.

%% ====================================================================
\section{Connection to Nelson's Program}
\label{sec:nelson}

Nelson~\cite{Nelson2015} observed that in PRA, the bounds in the
Hilbert--Ackermann consistency theorem depend on rank and level of
the proof, but \emph{not on the length of the proof}. For each
specific proof, the system can verify the relevant bounds. His planned
argument was to exploit the gap between this instance-by-instance
verification and a universal consistency assertion, using the
Kritchman--Raz ``surprise examination'' method~\cite{KritchmanRaz2010}
to derive a contradiction.

Our formalization makes this gap precise. Define:

\begin{definition}[Instance reflection]
For each proof code $p$ with $\checkPN(n{+}1, p) = \mathsf{just}(A)$,
the fuel-based system can produce a code~$q$ such that
$\checkPN(n{+}3, q) = \mathsf{just}(\fceq(\ccheck(\clit(p)),
\clit(\enc{A})))$.
\end{definition}

This is $\mathsf{represent\text{-}fuel}$, our D1 for the fuel
checker. It says: for each given proof code, the system can internally
certify what that code proves. Moreover, by D3, this certification can
itself be certified, iteratively with bounded fuel overhead.

\begin{definition}[Internal consistency]
$\mathsf{Con} = \fcAll(\fimp(\fceq(\ccheck(\mathsf{cvar}_0),
\clit(\enc{\fbot})),\, \fbot))$, asserting that no proof code proves
$\fbot$.
\end{definition}

The contrast:

\begin{theorem}[Instance--universal gap]
Instance reflection is available (for each specific~$p$), but internal
consistency $\mathsf{Con}$ cannot be derived from instance reflection
alone. The former uses \textsf{ax-eval} on closed code expressions;
the latter requires proving a fact about a universally quantified
code variable, which \textsf{ax-eval} cannot produce.
\end{theorem}

This is the formal content of Nelson's structural observation. The
passage from ``each given proof code can be individually checked'' to
``no proof code yields contradiction'' is precisely the G\"odel-II
step. In the stratified architecture, even instance reflection has
limits (the evaluator is blind to its own reflection codes). In the
fuel-based architecture, instance reflection works at all levels
including self-referentially (D3), but the universal consistency
statement still sits above the system.

The gap identified by Buss and Tao~\cite{Nelson2015} in Nelson's
later argument---that the disproof of certain stages ``seems to
require an assumption that $S^*$ is consistent''---corresponds in our
framework to the need for internal soundness when lifting from
instance verification to universal consistency. This is the barrier
that G\"odel's second theorem makes explicit: a system's ability to
verify each individual proof does not entail its ability to prove its
own global consistency.

\subsection{Open Consistency}

Nelson emphasized that the consistency of the \emph{open} (quantifier-free,
induction-free) fragment of PRA is provable by a simple valuation argument,
even though full consistency is not. Our formalization confirms this
precisely.

Define the open fragment $\mathsf{OpenProof}$ as the subsystem with only
$\mathsf{ax\text{-}refl}$, $\mathsf{ax\text{-}k}$, $\mathsf{ax\text{-}s}$,
and $\mathsf{mp}$---no~$\mathsf{gen}$, $\mathsf{cgen}$, or
$\mathsf{ax\text{-}eval}$. Define a propositional valuation:
\begin{align*}
  \mathsf{Good}(\fbot) &= \bot \\
  \mathsf{Good}(\mathsf{feq}(s,t)) &= \top \\
  \mathsf{Good}(\fimp(A,B)) &= \mathsf{Good}(A) \to \mathsf{Good}(B)
\end{align*}
with all other constructors mapped to~$\top$.

\begin{theorem}[Open soundness, \texttt{open-sound}]
If\/ $\mathsf{OpenProof}(A)$, then $\mathsf{Good}(A)$.
\end{theorem}

\begin{proof}
By structural induction.
$\mathsf{ax\text{-}refl}$: $\mathsf{Good}(\mathsf{feq}(t,t)) = \top$.
$\mathsf{ax\text{-}k}$: $\lambda x.\, \lambda \_.\, x$.
$\mathsf{ax\text{-}s}$: $\lambda f.\, \lambda g.\, \lambda a.\, f\;a\;(g\;a)$.
$\mathsf{mp}$: apply the induction hypotheses.
\end{proof}

\begin{corollary}[Open consistency, \texttt{open-consistent}]
$\mathsf{OpenProof}(\fbot) \to \bot$.
\end{corollary}

\begin{proof}
$\mathsf{open\text{-}sound}(\mathit{pf}) : \mathsf{Good}(\fbot) = \bot$.
\end{proof}

The contrast with the full system is sharp:
\begin{itemize}
\item The open fragment (propositional axioms + modus ponens) admits a
  four-line consistency proof via a propositional valuation.
\item The full system (with quantifier rules $\mathsf{gen}$, $\mathsf{cgen}$)
  does not admit an internal consistency proof (G\"odel~II).
\item The boundary is exactly the quantifier/induction rules: they cross
  the threshold where the simple valuation argument breaks, because
  quantified formulas require model-theoretic reasoning that the
  propositional valuation cannot provide.
\end{itemize}

This confirms Nelson's structural observation: there exists a meaningful
fragment weak enough to admit self-consistency, and the gap between this
fragment and the full system is precisely where the G\"odel barrier appears.

\subsection{The Corrected Nelson Program}

We package these results as a formal ``corrected Nelson theorem,''
preserving the valid content of Nelson's program while making the
barrier explicit.

\begin{theorem}[Corrected Nelson program, \texttt{ChwistekNelson.agda}]
\leavevmode
\begin{enumerate}[label=(\roman*)]
\item \emph{Open consistency.}
  $\mathsf{OpenProof}(\fbot) \to \bot$.
  The propositional fragment is consistent.

\item \emph{Instance-by-instance reflection.}
  For each code $p$ with $\checkPN(n{+}1, p) = \mathsf{just}(A)$,
  there exists a code $d$ with
  $\checkPN(n{+}3, d) = \mathsf{just}(\fceq(\ccheck(\clit(p)),
  \clit(\enc{A})))$.
  The system can certify what each specific code proves.

\item \emph{Iterated reflection.}
  D3 holds: the certification from~(ii) can itself be
  certified, at $+4$ fuel per level.

\item \emph{Meta-level consistency.}
  $\Proof(\fbot) \to \bot$.
  The full system is consistent (proved by soundness at the meta-level).

\item \emph{The gap.}
  Internal consistency ($\mathsf{Con}$ provable inside the system)
  is not derivable from instance reflection alone.
  $\mathsf{ax\text{-}eval}$ works on closed code expressions but cannot
  produce facts about universally quantified code variables.
  This is the G\"odel-II barrier.
\end{enumerate}
\end{theorem}

Nelson correctly identified that instance verification is available
and open consistency is provable. The gap between these and universal
internal consistency is real, precisely located, and is the formal
content of G\"odel's second incompleteness theorem.

%% ====================================================================
\section{Constructive G\"odel~I}
\label{sec:constructive}

The previous results show that $G$ is not provable (meta-level).
We now strengthen this to an \emph{internal proof transformation}:
given a proof of~$G$, \emph{construct a proof of}~$\fbot$.

This requires three additional rules in the proof system:
\begin{itemize}
\item $\mathsf{cinstC}$: from $\mathsf{ProofC}(\fcAll(A))$ and a code~$c$,
  derive $\mathsf{ProofC}(A[c])$ (universal elimination for code quantification).
\item $\mathsf{fceqTr}$: from $\fceq(e_1,e_2)$ and $\fceq(e_2,e_3)$,
  derive $\fceq(e_1,e_3)$ (transitivity).
\item $\mathsf{fceqSy}$: from $\fceq(e_1,e_2)$, derive $\fceq(e_2,e_1)$
  (symmetry).
\end{itemize}

\begin{theorem}[Constructive G\"odel~I, \texttt{constructive-goedel1}]
Given $\mathsf{ProofC}(G)$ and a code~$p$ with
$\checkPN(n{+}1,p) = \mathsf{just}(G)$, there is a
$\mathsf{ProofC}(\fbot)$.
\end{theorem}

\begin{proof}
Six steps:
\begin{enumerate}
\item $\mathsf{cinstC}$: instantiate $G = \fcAll(\mathit{body})$ at~$p$,
  obtaining $\fimp(\fceq(\ccheck(\clit(p)), \csub(\ldots)), \fbot)$.
\item $\mathsf{axEvalC}$: $\ccheck(\clit(p))$ evaluates to $\enc{G}$
  (from the encoding hypothesis).
\item $\mathsf{axEvalC}$: $\csub(\clit(\overline\varphi),
  \clit(\overline\varphi))$ evaluates to $\enc{G}$
  (self-reference, \texttt{selfRefGen}).
\item $\mathsf{fceqSy}$ on step~3: swap sides.
\item $\mathsf{fceqTr}$ on steps 2 and~4: chain the equalities.
\item $\mathsf{mpC}$ on steps 1 and~5: $\fbot$.
\end{enumerate}
\end{proof}

This is strictly stronger than the previous meta-level results:
\begin{center}
\begin{tabular}{@{}ll@{}}
\texttt{goedel1-final} & $\Proof(G) \to \bot$ \quad (meta-level) \\
\texttt{goedel1-fuel} & $\ProofN(G) \to \bot$ \quad (meta-level, fuel) \\
\texttt{constructive-goedel1} &
  $\mathsf{ProofC}(G) \to \mathsf{ProofC}(\fbot)$
  \quad (\textbf{internal})
\end{tabular}
\end{center}

The third produces an actual \emph{proof object of~$\fbot$} inside the
system, not just a meta-level contradiction. For full G\"odel~II
($\neg\,\mathsf{ProofC}(\mathsf{Con})$), one would additionally need
the internal implication $\mathsf{Con} \to G$, which requires universal
reasoning about all proof codes---the precise barrier identified in
Section~\ref{sec:nelson}.

%% ====================================================================
\section{Conclusion}
\label{sec:conclusion}

We have presented a fully constructive, self-contained formalization of
G\"odel's first incompleteness theorem in Agda, using binary trees as
the native syntax representation and avoiding G\"odel numbering entirely.
The development comprises 20 Agda files totaling approximately 4000
lines, with no postulates, no holes, and no standard library.

The key technical insight is that self-reference can be achieved through
a finite quine construction ($\csub$) operating on tree codes, without
any detour through arithmetic. The formalization revealed a sharp
structural dichotomy between stratified and fuel-based architectures:
the stratified approach gives an unconditional G\"odel~I and a strict
reflection hierarchy, while the fuel-based approach enables all three
derivability conditions and a bounded G\"odel~II.

The reflection hierarchy theorem---showing that each proof system can
reflect the layer below but is provably blind to its own reflective
reasoning---provides formal content to Chwistek's insistence on the
necessity of metasystem/object-system separation. In the fuel-based
architecture, this hierarchy collapses, aligning with Guard's
single-system approach, but at the cost of a fuel bound and an encoding
correctness assumption.

The open consistency result provides a precise lower bound: the
propositional fragment \emph{does} admit a four-line consistency proof
via a valuation argument. The remaining gap for a fully internal
G\"odel~II is the same gap that exists in any sufficiently strong
system: the quantifier and induction rules cross the threshold where
the valuation argument breaks, and the system cannot prove its own
full soundness. This is not a limitation of the architecture but the
very content of G\"odel's second incompleteness theorem.

\bigskip

\noindent\textbf{Artifact.}
The complete Agda development is available as 20 source files. All files
compile with Agda's \texttt{--without-K --exact-split} flags. The
\texttt{--safe} flag is used throughout; no postulates or unsafe pragmas
appear anywhere in the development.

%% ====================================================================
\bibliographystyle{plain}
\begin{thebibliography}{10}

\bibitem{Chwistek1939}
L.~Chwistek.
\newblock \emph{The Limits of Science: Outline of Logic and of the
  Methodology of the Exact Sciences}.
\newblock Kegan Paul, Trench, Trubner \& Co., London, 1935/1948.
\newblock (Original Polish edition: \emph{Granice nauki}, 1935.)

\bibitem{Goedel1931}
K.~G\"odel.
\newblock \"Uber formal unentscheidbare {S\"atze} der {Principia
  Mathematica} und verwandter {Systeme}~{I}.
\newblock \emph{Monatshefte f\"ur Mathematik und Physik}, 38:173--198,
  1931.

\bibitem{Guard1969}
J.~R. Guard.
\newblock The equivalence of two formulations of {G\"odel's} theorem.
\newblock \emph{Notre Dame Journal of Formal Logic}, 10(3):235--244,
  1969.

\bibitem{KritchmanRaz2010}
S.~Kritchman and R.~Raz.
\newblock The surprise examination paradox and the second incompleteness
  theorem.
\newblock \emph{Notices of the AMS}, 57(11):1454--1458, 2010.

\bibitem{Leigh2010}
G.~E. Leigh and M.~Rathjen.
\newblock The {Friedman--Sheard} programme in intuitionistic logic.
\newblock \emph{Journal of Symbolic Logic}, 77(3):777--806, 2012.

\bibitem{Nelson2015}
E.~Nelson.
\newblock Inconsistency of {Primitive Recursive Arithmetic}.
\newblock \emph{arXiv:1509.09209}, 2015.
\newblock With an introduction by S.~J.~Nelson and an afterword by
  S.~Buss and T.~Tao.

\bibitem{OConnor2005}
R.~O'Connor.
\newblock Essential incompleteness of arithmetic verified by {Coq}.
\newblock In \emph{Theorem Proving in Higher Order Logics (TPHOLs 2005)},
  LNCS 3603, pages 245--260. Springer, 2005.

\bibitem{Paulson2015}
L.~C. Paulson.
\newblock A mechanised proof of {G\"odel's} incompleteness theorems
  using {Nominal Isabelle}.
\newblock \emph{Journal of Automated Reasoning}, 55(1):1--37, 2015.

\bibitem{Shankar1994}
N.~Shankar.
\newblock \emph{Metamathematics, Machines, and {G\"odel's} Proof}.
\newblock Cambridge Tracts in Theoretical Computer Science 38.
  Cambridge University Press, 1994.

\end{thebibliography}

\end{document}
